\documentclass{pma-mthesis}
\begin{document}
\maketitlepage{
    theme = {Метод стиснення моделей глибинного навчання в умовах обмежених обчислювальних ресурсів},
    group = {КМ-41мн},
    author = {Передерей Богдан Олександрович},
    supervisorTitles = {канд. тех. наук, доцент кафедри ПМА},
    supervisor = {Тавров Данило Юрійович},
    normConsultantTitles = {старший викладач},
    normConsultant = {Мальчиков Володимир Вікторович},
% за відсутності консультанта приберіть/закоментуйте наступні три рядки
    % consultantIntro = {Консультант з \dots}, % розділи, "технічна частина" тощо
    % consultantTitles = {регалії консультанта},
    % consultant = {ПІБ},
    reviewerTitles = {регалії рецензента},
    reviewer = {ПІБ рецензента},
    year = {2026}
}

% завдання на дипломну роботу (впишіть відповідні значення)
\begin{assignment}{Передерею Богдану Олександровичу}
        {Богдан ПЕРЕДЕРЕЙ}
        {Данило ТАВРОВ}
    % у першому пункті потрібно вказати тему дипломної роботи, прізвище, ім’я, по батькові (повністю), науковий ступінь, учене звання керівника роботи, а також реквізити наказу, яким їх затверджено.
    \item Тема роботи: \inquot{\theme}, керівник роботи \supervisor{}, \supervisorTitles{}, затверджені наказом по університету від \dateplaceholder{} №~\_\_\_\_-\_. % від \inquot{15} квітня 2026 р. № 1808-C.

    % у другому пункті потрібно вказати термін подання переплетеної роботи на кафедру.
    \item Термін подання студентом роботи: \dateplaceholder{} % \inquot{11} червня 2025 р.

    % у пункті «Вихідні дані до роботи» потрібно вказати кількісні/якісні показники щодо умов, засобів та методів, які характеризують спрямованість дослідження, конкретизують методику розв'язання задачі та проведення експерименту
    \item Об’єкт дослідження: методи стиснення, модифікації архітектурних компонентів та підвищення ефективності виконання операцій нейронних мереж для їхнього навчання та розгортання в умовах обмежених ресурсів. 
    Дослідження охоплює етапи від адаптивного тренування до оцінки впливу дискретизації та зміни архітектури шарів на загальну поведінку нейронних мереж.

    % у пункті «Зміст роботи» потрібно вказати завдання з окремих частин дипломної роботи. Формулювати ці завдання потрібно в наказовому способі, починаючи зі слів «розробити», «виконати», «проаналізувати», «провести» тощо
    \item Предмет дослідження: алгоритми квантування у процесі навчання (Quantization Aware Training, QAT) 
    та методи апроксимації градієнтів для перетворення ваг нейронних мереж із форматів високої точності у числах з рухомою комою (Floating Point, FP)  до низької розрядності у цілочисельних представленнях (INT8, INT4). 
    
    \item Перелік завдань, які потрібно розробити: 
        \begin{itemize}[label=--, leftmargin=1.25cm, labelsep=1em]
            \item Провести огляд та аналіз літератури щодо сучасних методів стиснення та алгоритмічної перебудови нейронних мереж
            для мінімізація ресурсних витрат, особливостей їх застосування та впровадження у процес тренування та інференсу.
            \item Сформувати математичну постановку задачі квантування параметрів нейронної мережі та дослідити методи апроксимації градієнтів 
            для забезпечення стабільності навчання у форматах низької розрядності.
            \item Розробити експериментальний програмний комплекс для дослідження квантування у процесі навчання (QAT) на прикладі архітектури 
            згортокової мережі MobileNetV3.
            \item Провести серію порівняльних експериментів для оцінки впливу різних конфігурацій QAT на якість прогнозів та обчислювальну складність 
            моделі при переході до цілочисельних форматів із низькою розрядність (8 біт на менше).
            \item На основі отриманих результатів розробити та верифікувати методичний підхід (рекомендації) щодо ефективного застосування алгоритмів
            квантування для стиснення нейронних мереж.
        \end{itemize}
    
    % у пункті «Перелік ілюстративного матеріалу» потрібно вказати назви ключових для дипломної роботи рисунків та слайдів презентації, із якою студент виступатиме під час захисту дипломної роботи
    \item Перелік ілюстративного матеріалу: таблиця порівяння методів, графіки та таблиці метрик моделей під процесу тренування, 
    графіки та таблиці метрик натренованих моделей, схеми шарів та архітектур нейронних мереж, математичні формули та алгоритми, що описують
    методи квантування. 

    % Костиль, але поки так
    \item Орієнтовний перелік публікацій:  
    
    Perederei B., Qureshi F. Single Hyperspectral Image Super-Resolution Utilizing
          Implicit Neural Representations [Text] // Proceedings of the 14th International Conference 
          on Pattern Recognition Applications and Methods - Volume 1: ICPRAM / INSTICC. ––
          [S. l.] : SciTePress. –– 2025. –– P. 628–635.
    % у пункті «Консультанти розділів роботи» потрібно в окремих рядках таблиці вказати назви розділів, а також прізвища, ініціали та посади консультантів цих розділів (за наявності, у протилежному випадку цей пункт можна пропустити). Варто зазначити, що консультант із нормоконтролю є консультантом усієї роботи, а не окремих розділів, тому в цьому пункті його зазначати непотрібно
    % \item
    % \begin{consultants}
    %     Назва розділу & Прізвище та ініціали, посада &&\\\hline
    %     \dots & \dots &&
    % \end{consultants}

    % у пункті «Дата видачі завдання» потрібно вказати дату, коли студенту було видано завдання на дипломну роботу.
    \item Дата видачі завдання: \inquot{2} вересня 2025 р.
    
    % Stopped there
    \vspace{1.5em}
    \begin{calenderplan}
        % № з/п & Назва етапу & Термін виконання & Примітка
        1. & Аналіз досліджуваної галузі & 28.09.2025 & \\\hline
        2. & Формування структури дисертації, пошук та вивчення літератури & 29.10.2025 & \\\hline
        3. & Підготовка теоретичного розділу магістерської дисертації & 23.11.2025 & \\\hline
        4. & Планування наукового дослідження та робота над другим розділом магістерської дисертації & 18.12.2025 & \\\hline
        5. & Проведення наукового дослідження, розробка необхідного програмного забезпечення, збір та аналіз отриманих результатів & 21.03.2026 & \\\hline
        6. & Продовження проведення наукового дослідження, робота над третім та четвертим розділом дисертації & 19.04.2026 & \\\hline
        7. & Закінчення роботи над основною частиною магістерської дисертації & 16.05.2026 & \\\hline
        8. & Технічне та стилістичне оформлення пояснювальної записки та супровідних матеріалів & 24.05.2026 & 
    \end{calenderplan}
\end{assignment}

% «Анотація» повинна стисло відображати загальну характеристику та основний зміст дипломної роботи. Обсяг «Анотації» повинен становити не менше за 650 знаків (із пробілами).
\annotation{Реферат}
% «Анотацію» потрібно виконувати згідно з такою структурою:
% - відомості про обсяг дипломної роботи (БЕЗ додатків), кількість рисунків, таблиць, додатків і бібліографічних найменувань за «Переліком посилань»
Магістерську дисертацію виконано на \totpages{аркуші}{аркушах}{аркушах}, вона містить \totapp{додаток}{додатки}{додатків} та перелік посилань на використані джерела з \totcit{найменування}{найменувань}{найменувань}. У роботі наведено \totfig{рисунок}{рисунки}{рисунків} та \tottab{таблицю}{таблиці}{таблиць}.

\textbf{Актуальність теми}

Актуальність теми зумовлена стрімким розвитком глибинного навчання та критичною потребою в оптимізації моделей для реальних застосувань. 
Основною проблемою є масштабованість моделей, оскільки сучасні архітектури стають надто громіздкими для масового 
впровадження. Через гостру нестачу обчислювальних ресурсів для інференсу великих мовних моделей (LLM), виникає висока потреба
в методах, які дозволять мінмізувати використання обчислювальних витрат та пам'яті при збереженні функціональних можливостей моделей.

Також застосування таких методів дає можливість запускати нейронні мережі на периферійних пристроях, 
створюючи рішення для випадків, де приватність даних, низька затримка та автономність є вирішальними. 

\textbf{Зв’язок роботи з науковими програмами, планами, темами}

Робота над магістерською дисертацією проводилась згідно з науковим планом дослідницьких робіт кафедри прикладної математики 
Національного технічного університету України «Київський політехнічний інститут імені Ігоря Сікорського», 
а також у рамках дослідницької програми «Responsible AI for Ukraine Research Program» (Факультет прикладних наук УКУ у співпраці з Center for Responsible AI при NYU Tandon School of Engineering). 

\textbf{Мета і задачі дослідження}
Метою роботи є розробка та дослідження методів ефективного стиснення нейронних мереж для їх тренування та розгортання 
в умовах обмежених обчислювальних ресурсів. Основним завданням є критична оцінка готовності наявних методів до використання у реальних технічних задачах. 
Дослідження спрямоване на пошук наукових прогалин щодо практичного застосування таких технік та їх теоретичне і емпіричне заповнення. 
Окрему увагу приділено аналізу компромісів, які виникають під час використання таких методів, а саме між точністю, швидкістю та обсягом використання пам'яті.

\textbf{Об’єктом дослідження} є 

% - мета роботи
% - методи, розглянуті в роботі, критерії, за якими здійснювався вибір методу(-ів) для роботи, метод(-и), вибраний(-і) для використання в роботі
% - основні результати, яких було досягнуто в процесі виконання роботи. При цьому потрібно використовувати безособові форми дієслова доконаного способу минулого часу, наприклад, «виконано», «реалізовано», «проведено», «спроектовано» тощо
% - відомості щодо апробації результатів дипломної роботи (виступи на наукових конференціях, публікації в наукових журналах, довідки про впровадження тощо) чи рекомендації щодо використання результатів
% - ключові слова (від 5 до 20 позицій)

\annotation{Annotation}
% «Abstract» (анотація англійською мовою) повинна бути перекладом анотації українською мовою. Обсяг «Abstract» повинен становити не менше за 650 знаків (із пробілами).
This thesis is submitted on \totpages{page}{pages}{pages}, it has \totapp{appendix}{appendices}{appendices} and \totcit{reference}{references}{references}. There are \totfig{figure}{figures}{figures} and \tottab{table}{tables}{tables} in this document.

\tableofcontents

% «Вступ» має відображати актуальність роботи на основі стислого огляду сучасного стану проблемної області.
\intro

% Першим розділом роботи *обов'язково* повинен бути розділ «Постановка задачі», який повинен містити мету роботи, завдання, які потрібно розв'язати для її досягнення, вимоги до розроблюваних у роботі математичних методів та програмних засобів.
\chapter{Постановка задачі}

% У наступних розділах дипломної роботи потрібно подати:
% - огляд актуальних теоретичних та практичних рішень за тематикою роботи, критерії вибору методу(-ів) для використання в роботі, порівняльний аналіз розглянутих рішень за відповідними критеріями
% - опис побудованої математичної моделі та методу(-ів) розв'язання поставленої задачі
% - опис програмних засобів розв'язання поставленої задачі, у т. ч. розроблених алгоритмів
% - опис випробування розроблених програмних засобів
\chapter{Другий розділ}
Посилання додаються таким чином~\cite{Aho1992,Laming1968,NML,icpram2025}.

Математична формула у тексті: $x+\frac{x}{y} - y^{10} + y^20$, порівняйте степінь в $\{\dots\}$ та без. Формула окремим блоком з нумерацією:
\begin{equation}
    a = \frac{\partial b}{\partial{} y}\times\left.\sum_{i=0}^{\infty}\frac{1}{x^i}\right| _{x\rightarrow7} + 280\sin^3(\gamma)\left(\frac{\tanh(130^\circ - 42\pi n)}{\cosh(\alpha - \beta^2 + \sigma)}+\int_0^5 \frac{x^3 dx}{x + x^5} \right),
\end{equation}
\begin{eqdesc}
    \item[$a$] --- щось;
    \item[$b$] --- вже щось інше;
    \item[\dots] --- \dots.
\end{eqdesc}

Тестова таблиця (див. табл~\ref{tab:testing}).
\begin{table}{|l|r|}{Назва таблиці.}{tab:testing}{\hline a & b\\\hline}
    c & d
\end{table}

% Кожний розділ роботи, окрім «Постановки задачі», повинен закінчуватися підрозділом «Висновки до розділу».
\section{Висновки до розділу}

\chapter{Ще один розділ}

% У «Висновках» потрібно викласти найважливіші результати, одержані в дипломній роботі, потрібно зазначити якісні та кількісні показники відповідних результатів.
\conclusions

% «Перелік посилань» потрібно оформлювати за вимогами, викладеними на сторінці «Вимоги до оформлення документації».
\bibliographystyle{tex/udstu2006}
\bibliography{tex/bibliography}

% У дипломній роботі, серед інших, повинні бути такі обов'язкові додатки:
\appendix

% першим додатком повинен бути додаток «Лістинги програм», до якого потрібно винести лістинги програмного коду реалізації описаних у роботі математичних моделей, методів та алгоритмів. Код графічного інтерфейсу користувача можна не наводити. Не менше 30% рядків лістингів повинні складати коментарі. Обсяг лістингів повинен складати 15–40 сторінок. Варто пам'ятати, що обсяг лістингів не впливає на обсяг дипломної роботи!
\chapter{Лістинги програм}
\lstinputlisting[language=Python, caption=main.py]{../code/main.py} %< таким чином можна включити текст файлу до лістингів; TODO \srcset

% Додаткову інформацію: доведення теорем, матеріали до роботи --- потрібно також подавати у додатках
\chapter{Додаткові матеріали}

% передостаннім додатком повинен бути додаток «Ілюстративний матеріал», до якого потрібно винести слайди презентації, із якою студент виступатиме під час захисту дипломної роботи
\chapter{Ілюстративний матеріал}

% останнім додатком, за наявності відповідних матеріалів, повинен бути додаток «Довідки про впровадження результатів роботи», до якого потрібно винести скан-копії довідок (актів) про впровадження результатів дипломної роботи. У випадку відсутності таких матеріалів цей додаток можна не наводити
\chapter{Довідки про впровадження результатів роботи}

\end{document}